\chapter{What this is}

This chapter states what this research paper records and the boundary of what it claims.

\section{The research paper}

The aim is narrow: check one specific claim against its source documents and record what the check
finds.

Four things are in it. First the corpus, and what each document is once opened, as against what its
filename promises. Then the state of the Navier--Stokes problem as of September 2026. Then a reading
of the published blowup result against the official problem statement and against the Lean
formalization published with it. Last, an assessment of where an arbitrary-precision apparatus could
be aimed among the remaining problems.

No problem is attacked. Nothing is measured, run, or begun. Leading a research paper with a method
claims the method reaches something; no claim of that kind appears here or is supported here.

\section{What is not claimed}

No priority is claimed. The crystallography research paper~\cite{src:Quigg-crystallography} says the
same of its own construction.

The reading of the literature is not done. For any of the five open problems in the corpus, somebody
has to work through the field behind it and set down what that field already knows. Until then,
these pages assert no new method and credit no old one.

A priority dispute is running in the field the Navier--Stokes chapter covers. It is not investigated
here, no position is taken, and no line in that chapter is evidence on either side. The chapter
reports what the documents say and does not ask who reached the result first.

\section{Method}

Five rules govern the text, and each limits what the later chapters may claim.

\textbf{Cite nothing unread.} A document is cited here only once it is opened. A fact that arrives by
report is marked as reported in the sentence that carries it.

\textbf{A secondhand fact is no second witness.} A fact that arrives secondhand is confirmed against
the document before it is written down.

\textbf{Gaps go in the sentence making the claim.} Never in a footnote or an appendix. The
Navier--Stokes chapter has a section on what is not verified, and marks the individual soft spots
where they fall.

\textbf{A claim that does not hold stays on the page, with what refutes it.} The derived null in the
toolkit chapter is kept beside the drawn null that replaces it.

\textbf{Draw the bar. Never derive it.} A threshold reasoned out of an assumed distribution omits
every variance source not enumerated, and omitted variance only ever pushes one way. The toolkit
chapter applies the rule to the recommended measurement.
